\section*{Thông tin đề tài}
\addcontentsline{toc}{section}{Thông tin đề tài}

\noindent\begin{tabular}{@{}p{4.5cm}p{10.5cm}@{}}
\textbf{Mã đề tài:} & HK261-DAGD1-133 \\[2pt]
\textbf{Học phần:} & CO4041 -- Đồ án chuyên ngành / Đồ án môn học Kỹ thuật Máy tính, HK261 \\[2pt]
\textbf{GVHD:} & ThS. Phạm Công Thái, TS. Lê Trọng Nhân \\[2pt]
\textbf{Sinh viên thực hiện:} & Lê Thanh Phú (2212581), Phan Trần Nguyên Phúc (2212643), Hoàng Ngô Thiên Phúc (2212612) \\[2pt]
\textbf{Ngày cập nhật:} & 17/09/2026 \\
\end{tabular}

\section*{Tóm tắt}
\addcontentsline{toc}{section}{Tóm tắt}

Đề tài xây dựng một mô hình microgrid quy mô nhỏ dùng năng lượng mặt trời, gồm nguồn quang điện, bộ điều khiển sạc, pin lưu trữ, bộ quản lý pin (BMS), inverter, tải và lớp giám sát/điều khiển qua IoT. Báo cáo trình bày cơ sở lý thuyết đúc kết từ tài liệu tham khảo, phương pháp thiết kế BMS 12~V và mô phỏng BMS cùng tấm pin trên Proteus, thuật toán cấu hình lại dàn pin khi bức xạ không đồng đều, kiến trúc tổng thể của hệ thống, cùng kế hoạch và phân công cho toàn bộ học phần.

Mục tiêu của báo cáo là đề xuất hướng thiết kế và lộ trình triển khai để nhận góp ý của giảng viên. Thông tin pin hiện có là 12~V, 15~Ah; loại pin, tải và mã thiết bị triển khai còn cần xác định. Nhóm dự kiến tích hợp mô-đun sạc và inverter, tập trung phần tự phát triển vào BMS, đo lường, firmware, điều phối nguồn và IoT. Báo cáo trình bày phương pháp và kế hoạch; schematic Proteus, kết quả mô phỏng, số liệu đo và BOM có giá là đầu ra của các mốc tiếp theo.
